Prijanjanje (engl. \emph{perching}) je biološki nadahnuta sposobnost bespilotnih letjelica (UAV) da se privremeno pričvrste ili usidre na strukturu iz svoje okoline, umjesto da neprestano lete ili slijeću na tlo. Takva sposobnost korisna je za produljenje trajanja misije bez punjenja baterije, prikupljanje uzoraka, nadzor i inspekciju te niz drugih zadataka u kojima letjelica ima koristi od mogućnosti da zaustavi let, zadrži položaj na strukturi i po potrebi nastavi letjeti.

Problem prijanjanja UAV-a dobro je poznat u istraživačkoj zajednici, a predložena su brojna rješenja, no nekoliko se izazova stalno ponavlja. Prvi je prostorna preciznost: točku prijanjanja potrebno je točno odrediti i neprekidno pratiti dok je letjelica u pokretu, bez gubitka cilja uslijed poremećaja ili prebacivanja na pogrešnu metu. Drugi je izazov upravljanja: izvođenje manevra prijanjanja zahtijeva veliku preciznost pri brzini kretanja, gdje se male pogreške u prilazu izravno odražavaju na uspješnost. Treći je izazov računalne prirode, budući da implementacije percepcije za prijanjanje UAV-a zahtijevaju izvođenje modela dubokog učenja (npr. YOLO ili RT-DETR) koji trebaju značajnu računalnu snagu za rad u stvarnom vremenu usporedno s ostalim dijelovima obrade i upravljanja letjelicom.

Ovaj rad nastavak je i prilagodba autorova prethodnog diplomskog rada \cite{diplomski_rad}, u kojem je razvijen cjeloviti sustav za percepciju i upravljanje koji cijelim tokom obrade radi na samoj letjelici. U diplomskom radu granu je detektirao model dubokog učenja za segmentaciju instanci pokrenut na rubnom akceleratoru, dobiveni kandidat za prijanjanje pratio se metodom optičkog toka, a upravljanje prilazom izvodilo se vizualnim serviranjem temeljenim na slici (IBVS), uz konačni manevar prijanjanja okinut očitanjem senzora udaljenosti (ToF) te uz oslanjanje na sustav OptiTrack za apsolutnu lokalizaciju letjelice.

Natjecateljski scenarij kojem je namijenjen ovaj rad razlikuje se od scenarija diplomskog rada u nekoliko bitnih pretpostavki. Letjelica polazi već pozicionirana neposredno ispod grane, čime se ukida potreba za fazom vodoravnog prilaska grani. Sustav vanjske lokalizacije OptiTrack nije dostupan, pa se upravljanje mora osloniti isključivo na podatke dostupne na samoj letjelici. Senzor udaljenosti (ToF) ne koristi se za okidanje manevra prijanjanja, budući da je namijenjen drugoj svrsi izvan opsega ovog rada. Konačno, kamera je fizički postavljena da gleda prema gore, a ne prema naprijed kao u izvornoj konfiguraciji, čime se mijenja geometrija cijelog problema poravnanja i prilaska.

Ovaj rad opisuje ponovnu upotrebu i prilagodbu postojećeg cjevovoda za percepciju (detekcija grane, praćenje značajki) uz zamjenu upravljačkog sloja letjelice. Umjesto ROS čvora za upravljanje misijom prijanjanja razvijenog u sklopu diplomskog rada, koji je bio izravno ovisan o sustavu OptiTrack i o horizontalnoj geometriji prilaska, koristi se paket \texttt{ibvs\_perching}, otvoreni ROS paket koji izravno upravlja letjelicom preko MAVROS-a slanjem naredbi kutne brzine tijela i brzine penjanja, bez oslanjanja na vanjski sustav pozicioniranja ili na kontroler pozicije unutar leta.

\section{Povezani radovi}
\label{sec:povezani-radovi}

Detekcija grana za primjene na UAV-u dosad je istraživana u nekoliko smjerova. Silva i sur.~\cite{SILVA2022102759} predložili su metodu detekcije linija temeljenu na konvolucijskoj neuronskoj mreži za prepoznavanje grana stabla iz digitalnih fotografija, koristeći Houghovu transformaciju za procjenu orijentacije grane i točke prihvata. Lin i sur.~\cite{lin2025performanceevaluationdeeplearning} usporedili su nekoliko modela dubokog učenja za segmentaciju u autonomnim šumarskim sustavima, dajući usporedbu performansi modela u stvarnim uvjetima. Oba rada usredotočena su isključivo na percepcijsku stranu problema, bez obrade cjelokupnog cjevovoda od detekcije do upravljanja.

Na strani upravljanja, Ji i sur.~\cite{ji2022realtimetrajectoryplanningaerial} predstavili su metodu planiranja trajektorije u stvarnom vremenu za zračno prijanjanje. Samadikhoshkho i Lipsett~\cite{drones8110605} kombinirali su duboko učenje s vizualnim serviranjem temeljenim na slici za zračno uzorkovanje vegetacije. Domišlović i sur.~\cite{domislovic_marsupial} demonstrirali su vizualno serviranje temeljeno na poziciji s marsupijalnim robotskim sustavom -- pristup srodan onome korištenom u ovom radu. Kim i Chang~\cite{drones10040270} predstavili su sustav prijanjanja bez markera za kvadrokopter koji kombinira model segmentacije sa ToF senzorom i upravljačem temeljenim na konačnom automatu, izvodeći cijeli cjevovod na Jetson Orin NX platformi.

Za razliku od navedenih radova, koji redom pretpostavljaju letjelicu koja horizontalno prilazi meti uz raspoloživ sustav apsolutne lokalizacije, ovaj rad razmatra scenarij u kojem letjelica polazi već pozicionirana ispod cilja te se prijanjanje izvodi pretežno u vertikalnoj osi, bez oslanjanja na sustav vanjske lokalizacije poput OptiTracka.

\section{Pregled sustava}
\label{sec:pregled-sustava}

% TODO: umetnuti dijagram arhitekture sustava (analogno Slici 1 u diplomskom radu),
% prilagoden novoj konfiguraciji (kamera okrenuta prema gore, bez OptiTracka, bez ToF-a za prijanjanje)

Sustav zadržava podjelu na dva izvedbena okruženja iz diplomskog rada: Raspberry Pi OS, na kojem se izvodi percepcijski cjevovod, i Ubuntu unutar Docker kontejnera, na kojem se izvodi ROS upravljački sloj. Ta se podjela zadržava jer AI Hat+ akcelerator komunicira isključivo s Raspberry Pi OS okruženjem, pa sve što ovisi o slikovnim podacima ostaje na toj strani.

Na strani Raspberry Pi OS-a, cjevovod za detekciju grane preuzima okvire s kamere okrenute prema gore, provodi ih kroz model za segmentaciju instanci pokrenut na Hailo-8 akceleratoru te ih obrađuje kroz cjevovod naknadne obrade i ocjenjivanja kandidata kako bi se odredila jedinstvena točka prijanjanja. Ta se točka potom prosljeđuje IBVS algoritmu koji izračunava vektor pogreške u odnosu na središte slike.

Na strani upravljanja letjelicom, umjesto prethodnog ROS čvora za misiju prijanjanja koji je koristio OptiTrack odometriju i dvokoračnu trajektoriju (poravnanje -- vodoravni prilazak -- prijanjanje), koristi se paket \texttt{ibvs\_perching}. Taj paket izravno upravlja letjelicom preko sučelja \texttt{mavros/setpoint\_raw/attitude}, šaljući naredbe kutne brzine tijela za bočno poravnanje i naredbu brzine penjanja za vertikalni prilazak, bez posredovanja kontrolera pozicije ili vanjskog sustava lokalizacije.
